% TOP_PROBLEM: {"problemId": "problem.deterministic-parallel-perfect-matching-in-nc", "canonicalTitle": "Perfect-Matching Search in Deterministic NC (General Graphs)", "releaseRank": 227, "primaryDomain": "theoretical_computer_science", "primaryDomainLabel": "Theoretical computer science", "displayStatus": "open_with_solved_subcases", "releaseStatus": "open_with_solved_subcases"}
% !TeX program = lualatex
% Definition and sources only; this file is not an LLM solution attempt.
\documentclass[11pt]{article}
\usepackage[margin=25mm]{geometry}
\usepackage{fontspec}
\setmainfont{Latin Modern Roman}
\usepackage{amsmath,amssymb,mathtools}
\usepackage{unicode-math}
\setmathfont{Latin Modern Math}
\usepackage{xurl,hyperref}
\hypersetup{colorlinks=true,urlcolor=blue}
\setcounter{secnumdepth}{0}
\setlength{\parindent}{0pt}
\setlength{\parskip}{5pt}
\setlength{\emergencystretch}{3em}
\begin{document}
\section*{\texorpdfstring{Perfect-Matching Search in Deterministic NC (General Graphs)}{Perfect-Matching Search in Deterministic NC (General Graphs)}}\label{227}
\subsection{Definitions and mathematical statement}
A perfect matching in a finite simple undirected graph is a set of disjoint edges covering every vertex. On the adjacency-matrix encoding of a labeled $n$-vertex graph, seek one logspace-uniform family of deterministic Boolean circuits that outputs a matching matrix, or a distinguished NO flag precisely when no perfect matching exists. With fan-in at most two for AND and OR, the bounds are
\[
 \operatorname{size}(C_n)\le Cn^c,\qquad
 \operatorname{depth}(C_n)\le C(\log(n+2))^k
\]
for constants independent of $G,n$. A logarithmic-space transducer must print $C_n$ from $1^n$, including its designated outputs. This is a multi-output search problem on general graphs, not just testing existence or treating bipartite graphs.
\par\smallskip\textit{Formulation and concept references:} \href{https://arxiv.org/pdf/1704.01929}{[S1]}.

\subsection{Short English statement}
Can perfect matchings be found deterministically in general graphs by uniform polynomial-size circuits of polylogarithmic depth?
\subsection{Sources}
\begin{itemize}
\item[S1]Svensson and Tarnawski, The Matching Problem in General Graphs is in Quasi-NC, introduction; FGT Section 2.1 defines Search-PM; Nanongkai–Scquizzato Section 2.2 fixes logspace-uniform NC.
\url{https://arxiv.org/pdf/1704.01929}.
\textit{Formulation source.}
\end{itemize}
\end{document}
